% == == == == == == == == == == == == == == == == == == == == == == == == == ==
    == == == == == == == == == == == ==
    = % Section 3 : Conceptual Block Diagram % == == == == == == == == == == ==
      == == == == == == == == == == == == == == == == == == == == == == == == ==
      == ==
    =

\section {
  Conceptual Block Diagram
}
\sectionauthors { MN }

\subsection{System Overview} %
    Reviewed by Jeremie A Conceptual Block Diagram(CBD) provides a high-level visual representation of a system's major functional components and their relationships. Unlike detailed schematics, the CBD abstracts implementation specifics to communicate the overall system architecture and data flow to stakeholders. This top-down approach enables early validation of system concepts before committing to detailed design decisions.

The conceptual block diagram in Figure \ref{fig:cbd} illustrates the main sequence of operations performed by the system. To begin the process, the small mobile platform robot is placed within the environment to be mapped. The robot then systematically navigates throughout the space using an autonomous control algorithm. As it moves, it generates a map with its onboard LiDAR sensor and transfers the resulting data over the intranet to a computer for user access.

\begin{figure}[H]
\centering
\fbox{
  \includegraphics[width = 0.95\textwidth] { figures / CBD.png }
}
\caption { System Conceptual Block Diagram }
\label {
fig:
  cbd
}
\end { figure }
